% Finds the package and the example diagrams, also when compiled directly in an editor:
\makeatletter\def\input@path{{../}{../../desktop_app_rust/samples/}}\makeatother
\documentclass[10pt,a4paper]{article}
\usepackage[a4paper,margin=2cm]{geometry}
\usepackage[T1]{fontenc}
\usepackage{lmodern}
\usepackage{erdiagram}
\usepackage[most]{tcolorbox}
\tcbuselibrary{listings}
\usepackage{booktabs,array,longtable}
\usepackage{hyperref}
\hypersetup{colorlinks, linkcolor=blue!50!black, urlcolor=blue!50!black}

\lstdefinestyle{erd}{
  language=[LaTeX]TeX, basicstyle=\ttfamily\footnotesize, columns=fullflexible,
  keepspaces, commentstyle=\color{gray}\itshape,
  texcsstyle=*\color{blue!60!black}, keywordstyle=\color{blue!60!black},
  moretexcs={entity,relationship,isa,canbe,partof,attribute,attributes,link,
    point,note,subdomain,grid,setup,inputerdiagram,ersetup},
  literate={→}{{$\rightarrow$}}1,
}
\lstset{style=erd}
\newtcblisting{example}{enhanced, listing side text, colback=white,
  colframe=black!25, boxrule=0.5pt, arc=1pt, lefthand width=0.48\linewidth,
  listing options={style=erd}, halign lower=center, valign lower=center,
  before skip=6pt, after skip=8pt}
\newtcblisting{wideexample}{enhanced, listing above text, colback=white,
  colframe=black!25, boxrule=0.5pt, arc=1pt, listing options={style=erd},
  halign lower=center, before skip=6pt, after skip=8pt}
\newcommand\cs[1]{\texttt{\textbackslash#1}}
\newcommand\key[1]{\texttt{#1}}
\newcommand\meta[1]{\textlangle\textit{#1}\textrangle}
\newcommand\opt[1]{\texttt{[}#1\texttt{]}}

\title{\textbf{erdiagram} --- ER diagrams for \LaTeX\\[2pt]
  \large with n-notation, min-max cardinalities or crow's foot symbols, version 3.0.2}
\author{Leunami52}
\date{2026-11-07}

\begin{document}
\maketitle
\begin{center}\small Distributed under the \LaTeX\ Project Public License, version 1.3c or later.\end{center}
\tableofcontents

% =====================================================================
\section{Introduction}

\texttt{erdiagram} draws entity--relationship diagrams in the classic
style (diamonds for relationship types, attribute lists or Chen's ovals),
extended by a few useful constructs. It is built on TikZ, but you never have to write TikZ
yourself: a diagram is a list of short, declarative commands inside an
\texttt{erdiagram} environment.

\paragraph{Installation.} Copy \texttt{erdiagram.sty} next to your
document and load it with \verb|\usepackage{erdiagram}|. The package needs
only TikZ, which is part of every \TeX\ distribution. It can be loaded
together with its sister package \texttt{gsmdiagram}.

\paragraph{No name clashes.} The package defines only four public
names: \cs{inputerdiagram}, \cs{ersetup} and the environments \texttt{erdiagram}
and \texttt{relationalschema}. The short diagram commands (\cs{entity}, \cs{link},
\cs{note}, \dots) exist \emph{only inside} the environment, so they never
interfere with the rest of your document.

\begin{example}
\begin{erdiagram}
  \entity{building}(0,0){Building}
  \relationship[dependency=Id]{has}(3,0){has}
  \entity[weak]{room}(6,0){Room}
  \link{building}{has}{1}
  \link[double]{has}{room}{n}
  \attributes[anchor=north]{building}
    {*buildingID, address}
  \attributes[anchor=north]{room}
    {*roomNo, size}
\end{erdiagram}
\end{example}

Every command follows the same pattern
\[
  \cs{command}\opt{\meta{options}}\{\meta{id}\}(\meta{x},\meta{y})\{\meta{label}\}
\]
where the \meta{id} is the internal name used to connect elements, the
optional position $(x,y)$ is a TikZ coordinate (in cm) and the
\meta{label} is what is printed. Options are \texttt{key=value} pairs;
options the package does not know are passed on to TikZ, so everything
TikZ offers is still available.

% =====================================================================
\section{New in version 3.0}

Diagrams from earlier versions keep compiling. These changes show in them:

\begin{itemize}
\item IS\_A and CAN\_BE are hexagons; \verb|\ersetup{isa shape=diamond}| brings back the
  diamonds (section~\ref{sec:generalisation}).
\item The aggregation is labelled PART\_OF in capitals, and keys of weak entity types are dashed
  unless they depend on E alone.
\end{itemize}

And these are new:

\begin{itemize}
\item Cardinalities: n-notation and min-max with a defined meaning, and crow's foot symbols
  (section~\ref{sec:links}). The older forms \texttt{m}, \texttt{mc} and \texttt{[a,b]} keep working.
\item Composite, multi-valued and derived attributes in lists, further candidate keys, and lists drawn as
  ovals (section~\ref{sec:attributes}).
\item The types of a generalisation as words (section~\ref{sec:generalisation}).
\item \texttt{sans serif}, a real \texttt{scale}, and labels that follow the language of the document
  (section~\ref{sec:settings}).
\item The environment \texttt{relationalschema} for relational schemas (section~\ref{sec:rm}).
\end{itemize}

\section{The environment}

\begin{center}
\verb|\begin{erdiagram}[|\meta{options}\verb|] ... \end{erdiagram}|
\end{center}

\begin{longtable}{@{}>{\ttfamily}p{3.3cm}p{12.5cm}@{}}
\toprule
\normalfont\textbf{Option} & \textbf{Meaning}\\\midrule
grid & draws a labelled coordinate grid behind the whole diagram (a
  drafting aid for finding positions)\\
width=\meta{dim} & scales the finished diagram to this width\\
height=\meta{dim} & scales the finished diagram to this height\\
max width=\meta{dim} & shrinks the diagram if it is wider (never enlarges)\\
max height=\meta{dim} & shrinks the diagram if it is higher\\
scale=\meta{factor} & scales the finished diagram with its nodes, text and lines
  (section~\ref{sec:settings})\\
checks=false & switches off the rule checks at the end of the diagram
  (Section~\ref{sec:rules})\\
\meta{style key} & any key of \cs{ersetup} (Section~\ref{sec:style}),
  valid for this diagram only\\
\meta{TikZ option} & anything else is a \texttt{tikzpicture} option,
  e.g.\ \texttt{baseline}, \texttt{node distance=2cm}\\
\bottomrule
\end{longtable}

Positions can be given as absolute coordinates \verb|(x,y)|, as TikZ
expressions like \verb|($(a)+(2,1)$)|, or with the TikZ positioning keys
as options, e.g.\ \verb|\entity[right=2cm of a]{b}{B}|. Inside the
environment, \cs{setup}\verb|{...}| changes style keys for the rest of
the diagram, and arbitrary TikZ commands (\cs{draw}, \cs{node}, \dots)
may be mixed in.

\begin{example}
\begin{erdiagram}[grid, scale=0.8,
                  card placement=beside]
  \entity{a}(0,0){A}
  \relationship{r}(3,1){R}
  \entity{b}(6,0){B}
  \link{a}{r}{1}
  \link{r}{b}{cn}
\end{erdiagram}
\end{example}

% =====================================================================
\section{Diagrams in external files}\label{sec:files}

\begin{center}
\verb|\inputerdiagram[|\meta{environment options}\verb|]{|\meta{file}\verb|}|
\end{center}

The file contains only the \emph{body} of an \texttt{erdiagram}
environment, i.e.\ the lines between \verb|\begin{erdiagram}| and
\verb|\end{erdiagram}|. This is the format written by the desktop
editor, so diagrams can be edited graphically and included in a document
without copying any code. The recommended file extension is
\texttt{.erd}. Paths are relative to the main document; with the
\texttt{subfiles} package use \verb|\inputerdiagram{\subfix{diagram.erd}}|.

\begin{lstlisting}
% file: university.erd -- body of an erdiagram environment
\entity{student}(0,0){Student}
\relationship{attends}(3.5,0){attends}
\entity{course}(7,0){Course}
\link{student}{attends}{cn}
\link{attends}{course}{n}
\end{lstlisting}

\begin{example}
\inputerdiagram[width=0.9\linewidth]
  {university.erd}
\end{example}

\paragraph{Rules for files that the editor should be able to read
back.} One command per line; positions as absolute coordinates
\verb|(x,y)|; options as \texttt{key=value} pairs (values containing
commas in braces); \% comments are allowed. Plain TikZ code works in a
file as well, but the editor will not be able to edit it graphically.

% =====================================================================
\section{Entities}

\begin{center}
\verb|\entity[|\meta{options}\verb|]{|\meta{id}\verb|}(|\meta{x,y}\verb|){|\meta{label}\verb|}|
\end{center}

\begin{longtable}{@{}>{\ttfamily}p{3.3cm}p{12.5cm}@{}}
\toprule
\normalfont\textbf{Option} & \textbf{Meaning}\\\midrule
weak & weak entity type (double border)\\
color=\meta{colour} & background colour. A plain colour name is used as a
  light tint (\texttt{blue} $\to$ \texttt{blue!20}); a full xcolor
  expression such as \texttt{orange!35} is used as is\\
\meta{TikZ option} & e.g.\ \texttt{minimum width=4cm}, \texttt{draw=red},
  \texttt{right=of x}\\
\bottomrule
\end{longtable}

\begin{example}
\begin{erdiagram}
  \entity{a}(0,0){Entity}
  \entity[weak]{b}(3.2,0){Weak entity}
  \entity[color=blue]{c}(0,-1.5){Tinted}
  \entity[weak, color=orange!35]{d}(3.2,-1.5)
     {Two\\lines}
  \entity[minimum width=1cm, draw=red]{e}
     (1.6,-3){Small}
\end{erdiagram}
\end{example}

% =====================================================================
\section{Relationships}

\begin{center}
\verb|\relationship[|\meta{options}\verb|]{|\meta{id}\verb|}(|\meta{x,y}\verb|){|\meta{label}\verb|}|
\end{center}

\begin{longtable}{@{}>{\ttfamily}p{3.3cm}p{12.5cm}@{}}
\toprule
\normalfont\textbf{Option} & \textbf{Meaning}\\\midrule
weak & weak relationship type (double border)\\
dependency=E & existence dependency: double border and a red
  \textbf{E} above the name. The weak entity's own key is unique everywhere, so it stays
  solidly underlined\\
dependency=Id & identification dependency, which includes existence dependency: double border
  and a red \textbf{Id} above the name. The weak entity's key is only unique together with its
  owner's, so it is underlined with dashes\\
color=\meta{colour} & background colour, as for entities\\
\bottomrule
\end{longtable}

\begin{example}
\begin{erdiagram}
  \relationship{r1}(0,0){owns}
  \relationship[weak]{r2}(2.5,0){has}
  \relationship[dependency=E]{r3}(0,-1.8){has}
  \relationship[dependency=Id]{r4}(2.5,-1.8)
     {has}
  \relationship[color=yellow]{r5}(1.25,-3.6)
     {is under\\treatment}
\end{erdiagram}
\end{example}

Relationships may have any arity; just link all participating entities:

\begin{example}
\begin{erdiagram}
  \entity{test}(0,1.5){Test}
  \entity{doctor}(5,1.5){Doctor}
  \entity{patient}(2.5,-1.8){Patient}
  \relationship{reg}(2.5,0){registered}
  \link{test}{reg}{n}
  \link{doctor}{reg}{1}
  \link{patient}{reg}{n}
  \attributes{reg}{date}
\end{erdiagram}
\end{example}

% =====================================================================
\section{Links and cardinalities}\label{sec:links}

\begin{center}
\verb|\link[|\meta{options}\verb|]{|\meta{from}\verb|}{|\meta{to}\verb|}{|\meta{cardinality}\verb|}|
\end{center}

The cardinality argument is optional. By default it sits \emph{on} the
middle of the edge, on a white background. Any text is printed as given, but two
notations have a defined meaning, also for the relational schema of the desktop
editor (section~\ref{sec:rm}):

\begin{description}
\item[n-notation] \texttt{1} (exactly one), \texttt{c} (at most one), \texttt{n}
  (one or more), \texttt{cn} (any number). The label says how many entities of
  the type \emph{at its end of the edge} one partner has: ``Halter \texttt{1} --
  \texttt{n} Hund'' means each Hund has exactly one Halter.
\item[min-max notation] $(\min,\max)$, such as \texttt{(1,*)} or \texttt{(0,1)}: how
  often an entity \emph{of the type at its end} takes part. So the same fact is
  written ``Halter \texttt{(1,*)} -- \texttt{(1,1)} Hund'': the values sit at the opposite ends.
  \texttt{(·,*)} stands for an unknown minimum.
\end{description}

For relationships of a higher degree the two are not just flipped: \texttt{1} says
that the combination of all other sides determines this one, while $(\min,\max)$ only
counts how often one entity takes part. Older forms (\texttt{m}, \texttt{mc},
\texttt{[3,4]}) keep working.

\paragraph{Crow's foot symbols.} With \texttt{card symbols=crowsfoot} (an environment
option, or \cs{setup}, or \cs{ersetup} for all diagrams), n-notation is drawn as symbols at the
entity end of the edge, with the same meaning: \texttt{1} is two bars, \texttt{c} a circle and
a bar, \texttt{n} a bar and a fork, \texttt{cn} a circle and a fork. Only the labels are replaced:
diamonds and attributes stay, so relationships of a higher degree keep working. The edge
meets the side of the entity at a right angle: it ends in a short straight piece, which carries
the symbol, and the rest of the edge leads to that. All other labels stay text.

\begin{example}
\begin{erdiagram}[card symbols=crowsfoot]
  \entity[minimum width=1.6cm]{a}(0,0){Halter}
  \relationship{r}(2.7,0){besitzt}
  \entity[minimum width=1.6cm]{b}(5.4,0){Hund}
  \link{a}{r}{1}
  \link{r}{b}{n}
  \entity[minimum width=1.6cm]{c}(0,-2){Labor}
  \relationship{s}(2.7,-2){leitet}
  \entity[minimum width=1.6cm]{d}(5.4,-2){Forscher}
  \link{c}{s}{c}
  \link{s}{d}{cn}
\end{erdiagram}
\end{example}

\begin{example}
\begin{erdiagram}[card symbols=crowsfoot]
  \entity[minimum width=1.6cm]{a}(0,0){Halter}
  \relationship{r}(2.7,1.5){besitzt}
  \entity[minimum width=1.6cm]{b}(5.4,0){Hund}
  \link{a}{r}{1}
  \link{r}{b}{n}
\end{erdiagram}
\end{example}

\paragraph{Cardinalities beside the edge.} With
\texttt{card placement=beside} -- as an environment option for one
diagram, with \verb|\ersetup{card placement=beside}| for all diagrams,
or as the link option \texttt{beside} for a single edge -- the
cardinality is printed next to the \emph{entity} end of the edge, no matter in which order \meta{from} and \meta{to}
are given. With \texttt{card side=auto} it is placed on the side of its
edge with more free room: the package looks at all edges and attribute
trees leaving the same entity. If there is plenty of room on both sides
it prefers the position \emph{above} a flat edge and \emph{right} of a
steep one. The label never crosses its edge and is pushed outwards along
the edge until it does not overlap its entity.

\begin{example}
\begin{erdiagram}[card placement=beside,
                  max width=\linewidth]
  \entity{p}(0,0){Person}
  \relationship{lives}(3,1.3){lives}
  \relationship{born}(3,-1.3){born}
  \entity{c}(6,0){City}
  \link{p}{lives}{n}
  \link{lives}{c}{n}
  \link{p}{born}{n}
  \link{born}{c}{1}
  \attributes[anchor=north, xshift=-0.6cm,
    side=left]{p}{name, insurance number}
  \attributes[anchor=north east]{c}
    {zip code, name, population}
  \attributes[anchor=south]{born}
    {date of birth}
\end{erdiagram}
\end{example}

\begin{longtable}{@{}>{\ttfamily}p{3.6cm}p{12.2cm}@{}}
\toprule
\normalfont\textbf{Option} & \textbf{Meaning}\\\midrule
beside & this edge's cardinality next to the entity
  (\texttt{card placement=beside} does this for a whole diagram)\\
inline & this edge's cardinality on the middle of the edge (the default)\\
card side=\meta{side} & beside placement: \texttt{auto} (default),
  \texttt{above}, \texttt{below}, \texttt{left} or \texttt{right}\\
card dist=\meta{dim} & beside placement: distance from the entity border
  (default 0.3cm)\\
card pos=\meta{fraction} & position along the edge (0 = start, 1 = end),
  overrides \texttt{card dist}\\
role=\meta{text} & role name (needed for recursive relationships)\\
role side=\meta{side}, role pos & side and position (fraction) of the
  role name\\
bend=\meta{angle} & curved edge; positive bends to the left, negative to
  the right (seen in the direction of the edge)\\
via=\{(x,y) (x,y) \dots\} & route through way points\\
ortho=-| \textrm{or} |- & right-angled edge (horizontal-then-vertical or
  vice versa)\\
arrow & draw an arrow head at the \meta{to} end\\
double & double line. Good practice for the link between a weak
  relationship (weak, E or Id) and its weak entity, so that the chain of
  dependencies is visible. If a weak relationship joins two weak entities, only the link to
  the dependent one is double, which is how the desktop editor knows who depends on whom.
  The package warns if a double line joins
  anything else\\
\meta{TikZ option} & e.g.\ \texttt{bend left=30}, \texttt{out=90, in=180},
  \texttt{looseness=1.5}, \texttt{dashed}, \texttt{thick}\\
\bottomrule
\end{longtable}

\begin{example}
\begin{erdiagram}
  \entity{e}(0,0){Employee}
  \relationship{sup}(4,0){is super-\\visor of}
  \link[role=supervisor, bend=25]{e}{sup}{1}
  \link[role=employee, bend=-25,
        role side=below]{e}{sup}{n}
\end{erdiagram}
\end{example}

\begin{example}
\begin{erdiagram}
  \entity{a}(0,0){A}
  \relationship{r}(4,1.8){R}
  \entity{b}(4,-1.5){B}
  \link[ortho=|-, beside]{a}{r}{(0,n)}
  \link[via={(6,1.8) (6,-1.5)}]{r}{b}{(1,1)}
  \relationship{s}(0,-2.5){S}
  \link{a}{s}{c}
  \link[bend=-30, beside]{s}{b}{cn}
\end{erdiagram}
\end{example}

Edges can start at a particular anchor of a node
(\verb|\link{a.north east}{r}|), and \cs{point}\verb|{|\meta{id}\verb|}(x,y)|
creates an invisible junction to which several edges can be routed.

% =====================================================================
\section{Attributes}\label{sec:attributes}

\subsection{Attribute trees}

\begin{center}
\verb|\attributes[|\meta{options}\verb|]{|\meta{node}\verb|}{|\meta{attribute}\verb|, |\meta{attribute}\verb|, ...}|
\end{center}

The list is printed top to bottom. A leading \texttt{*} marks a key
attribute (underlined). Put an attribute in braces if it contains a
comma; any formatting is allowed.

\begin{longtable}{@{}>{\ttfamily}p{3.6cm}p{12.2cm}@{}}
\toprule
\normalfont\textbf{Option} & \textbf{Meaning}\\\midrule
anchor=\meta{anchor} & where the tree leaves the node: any TikZ anchor
  (\texttt{north}, \texttt{south east}, \dots) or an angle
  (\texttt{60}). Default: \texttt{east}\\
dir=\meta{direction} & \texttt{up}/\texttt{down}: vertical trunk with
  ticks; \texttt{left}/\texttt{right}: short stem and a vertical bar.
  Default: taken from the anchor\\
side=\meta{side} & trunk layout (\texttt{left} or \texttt{right}): on which side the labels are. Default:
  left for anchors containing \texttt{west}, right otherwise\\
align=\meta{where} & bar layout: where the stem meets the bar
  (\texttt{center}, \texttt{top} or \texttt{bottom})\\
spread, gap, tick, stem & distances between attributes, from the node to
  the first attribute, tick length, stem length\\
xshift, yshift & move the start of the tree\\
\meta{TikZ option} & applied to the attribute labels, e.g.\
  \texttt{font=\textbackslash scriptsize}, \texttt{text=blue}\\
\bottomrule
\end{longtable}

\begin{example}
\begin{erdiagram}
  \entity{e}(0,0){Entity}
  \attributes[anchor=north]{e}{*north, up}
  \attributes[anchor=south]{e}{south, down}
  \attributes[anchor=east]{e}{east, bar, right}
  \attributes[anchor=west]{e}{west, bar, left}
  \attributes[anchor=north east]{e}{NE}
  \attributes[anchor=north west]{e}{NW}
  \attributes[anchor=south east]{e}{SE}
  \attributes[anchor=south west]{e}{SW}
\end{erdiagram}
\end{example}

\begin{example}
\begin{erdiagram}
  \entity{e}(0,0){Employee}
  \attributes[anchor=north, xshift=-0.9cm,
    side=left]{e}{*empNo, name}
  \attributes[anchor=south, xshift=-0.9cm]{e}
    {{job, current}, \textit{salary}}
  \attributes[anchor=east, align=top]{e}
    {top, aligned}
  \relationship{r}(4,0){works}
  \attributes[anchor=north]{r}{since}
  \attributes[anchor=-60, text=blue]{r}
    {at -60 degrees}
\end{erdiagram}
\end{example}

\subsection{Composite, multi-valued and derived attributes}

Lists can hold all kinds of attributes. Parts follow their attribute in braces;
a multi-valued attribute is written \verb|\{email\}| and a derived one \verb|/age|. A
derived attribute can also say how it is computed (\verb|/age = {2026 - born}|), which only
the relational schema of the desktop editor uses. Other braces keep their meaning:
\verb|\textit{x}|, \verb|$x_{i}$| and \verb|ratio 1{:}2| are one name each.

\begin{example}
\begin{erdiagram}
  \entity{p}(0,0){Person}
  \attributes{p}{*PID, address{street, town},
    \{email\}, /age = {2026 - born}, born}
\end{erdiagram}
\end{example}

\subsection{Further candidate keys}

An entity type can have several candidate keys, of which one is drawn as its key.
Mark the others with a plus: \verb|+email|. Nothing of this is drawn; the desktop
editor makes them \texttt{UNIQUE} constraints. A name that starts with a number, such as \verb|+49 phone|,
stays a name. In an oval, it's the option \texttt{unique}.

\subsection{Keys of weak entity types}

The key of a weak entity type is only unique together with its owner, and is
drawn with a dashed underline, without anything to set: for the attributes of an
\cs{entity}\texttt{[weak]}, in a list as well as in an oval. It stays solid if the entity
depends on \texttt{dependency=E} alone, since that key is unique by itself. If an entity has
double-lined links, only those say what it depends on, so being the owner of another weak entity
over a single line doesn't change its key. The line is drawn at the end of the diagram, when all
links are known.

\begin{example}
\begin{erdiagram}
  \relationship[dependency=Id]{r}(0,0){has}
  \entity[weak]{a}(2.8,0){Room}
  \attributes[anchor=north]{a}{*number, seats}
  \link[double]{r}{a}{n}
  \relationship[dependency=E]{s}(0,-2.6){has}
  \entity[weak]{b}(2.8,-2.6){Copy}
  \attributes[anchor=north]{b}{*number, seats}
  \link[double]{s}{b}{n}
\end{erdiagram}
\end{example}

\subsection{Lists as ovals}
\label{sec:chenlists}

\texttt{attribute style=chen} draws every list as ovals, in Chen's notation, from the
same file. They are placed at the end of the diagram, in the directions around the node
that are furthest from its edges, clockwise from the top in the order of the list.
A composite attribute is a rounded box with its parts fanned out behind it. To fix an oval,
add its position in cm after the name, relative to what it belongs to:
\verb|name@2/1|. The desktop editor does that when an oval is dragged. Compact lists
ignore these positions, so a file can go back and forth.

\begin{example}
\begin{erdiagram}[attribute style=chen]
  \entity{p}(0,0){Person}
  \attributes{p}{*PID, name, address{street, town},
    \{email\}, /age}
\end{erdiagram}
\end{example}

\subsection{Chen-style attributes}

\begin{center}
\verb|\attribute[|\meta{options}\verb|]{|\meta{id}\verb|}(|\meta{x,y}\verb|){|\meta{label}\verb|}|
\end{center}

The classic notation (Chen's): attributes as
ellipses. Options: \texttt{key} (underlined), \texttt{unique} (a further
candidate key, not drawn), \texttt{derived} (dashed), \texttt{formula=}\meta{SQL} (how a
derived attribute is computed, only used by the relational schema of the desktop
editor), \texttt{multivalued} (double border), \texttt{composite} (rounded box for
aggregated attributes), \texttt{type=}\meta{type} (section~\ref{sec:types}) and
\texttt{of=}\meta{node}, which draws the connecting line.

\begin{example}
\begin{erdiagram}
  \entity{car}(0,0){Passenger car}
  \attribute[key, of=car]{a1}(-2.2,1.5){carNo}
  \attribute[of=car]{a2}(0,1.7){color}
  \attribute[derived, of=car]{a3}(2.2,1.5){age}
  \attribute[multivalued, of=car]{a4}
    (-2.4,-1.4){owner}
  \attribute[composite, of=car]{a5}
    (1.8,-1.5){engine}
  \attribute[of=a5]{a6}(1,-2.9){power}
  \attribute[of=a5]{a7}(3,-2.9){type}
\end{erdiagram}
\end{example}

\subsection{Data types}\label{sec:types}

Attributes can carry a data type, used by the desktop editor for the
relational schema and SQL export (section~\ref{sec:rm}). In an attribute
list the type follows a colon: \verb|*PID : integer|; an ellipse takes the
option \texttt{type=}\meta{type}. Types are \emph{not} printed unless
\texttt{show types} is set (as environment option, with \cs{ersetup} or
with \cs{setup}). A colon that belongs to the name must be braced
(\verb|ratio 1{:}2|), as must a type containing a comma
(\verb|price : {decimal(8,2)}|).

\begin{example}
\begin{erdiagram}[show types]
  \entity{p}(0,0){Person}
  \attributes[anchor=north]{p}
    {*PID : integer, name : string,
     born : date}
  \attribute[type=float, of=p]
    {h}(3.2,0){height}
\end{erdiagram}
\end{example}

\texttt{type sep} (default \verb*|:\ |) is put between name and type,
\texttt{type font} (default empty) selects the font of the type.

% =====================================================================
\section{Generalisation and aggregation}\label{sec:generalisation}

\begin{center}
\verb|\isa[|\meta{options}\verb|]{|\meta{id}\verb|}(|\meta{x,y}\verb|){|\meta{supertype}\verb|}{|\meta{constraint}\verb|}|\qquad
\verb|\canbe[...]{...}(...){...}{...}|
\end{center}

\cs{isa} draws the IS\_A hexagon and the arrow to the supertype;
subtypes are then connected with \cs{link}. \cs{canbe} is the same for
specialisations (CAN\_BE). With \texttt{isa shape=diamond} both are drawn as diamonds.

The constraint gives the type of the generalisation: \texttt{c} complete (also called
total) or \texttt{i} incomplete (also called partial), and \texttt{d} disjoint or
\texttt{o} overlapping, so \texttt{cd}, \texttt{co}, \texttt{id}, \texttt{io}. By
default it is printed inside the shape, below IS\_A. With
\texttt{constraint pos=auto} it goes beside the shape, in the direction
with the most room, or to a fixed side such as
\texttt{constraint pos=east}. \texttt{constraint style=words} prints
``complete, disjoint'' instead of the letters. Words go beside the shape; other
text is printed as written.

\begin{longtable}{@{}>{\ttfamily}p{3.6cm}p{12.2cm}@{}}
\toprule
\normalfont\textbf{Option} & \textbf{Meaning}\\\midrule
constraint pos=\meta{pos} & \texttt{inside} (default), \texttt{auto}
  or a compass direction like \texttt{east}\\
super card=\meta{text} & cardinality at the supertype end of the arrow
  (alternative notation: \texttt{1})\\
sub card=\meta{text} & cardinality below the shape (\texttt{n},
  \texttt{c}, \texttt{cn})\\
\meta{other} & as for relationships (\texttt{color}, TikZ options)\\
\bottomrule
\end{longtable}

\begin{example}
\begin{erdiagram}
  \entity{p}(0,0){Product}
  \isa{i}(0,-1.8){p}{cd}
  \entity{f}(-1.6,-3.5){Foreign material}
  \entity{s}(1.6,-3.5){Self produced}
  \link{f}{i}\link{s}{i}
\end{erdiagram}
\end{example}

The alternative notation with cardinalities instead of letters:

\begin{example}
\begin{erdiagram}
  \entity{q}(0,0){Player}
  \isa[super card=1, sub card=m]{j}(0,-1.8){q}
  \entity{x}(-1.6,-3.5){Soccer}
  \entity{y}(1.6,-3.5){Tennis}
  \link{x}{j}\link{y}{j}
\end{erdiagram}
\end{example}

\begin{example}
\begin{erdiagram}
  \entity{p}(0,0){Person}
  \isa[constraint pos=auto]{i1}(-2,-1.6){p}{cd}
  \isa{i2}(2,-1.6){p}{io}
  \entity[minimum width=1.6cm]{t}(-3,-3.2){Teen}
  \entity[minimum width=1.6cm]{a}(-1,-3.2){Adult}
  \entity[minimum width=1.6cm]{w}(1.1,-3.2){Worker}
  \entity[minimum width=1.6cm]{o}(3,-3.2){Officer}
  \link{t}{i1}\link{a}{i1}\link{w}{i2}\link{o}{i2}
\end{erdiagram}
\end{example}

\paragraph{Several supertypes.} A subtype can belong to several generalisations: draw one
\cs{isa} per supertype and link the subtype to each. What this means for the relational schema
is described in section~\ref{sec:rm}.

\begin{center}
\verb|\partof[|\meta{options}\verb|]{|\meta{id}\verb|}(|\meta{x,y}\verb|)|
\end{center}

An aggregation diamond labelled PART\_OF. \verb|\ersetup{partof label=part of}| changes the
text. It stays a diamond whatever \texttt{isa shape} says.

\begin{example}
\begin{erdiagram}
  \entity{car}(1.6,0){Passenger car}
  \partof{p1}(0,-1.6)\partof{p2}(3.2,-1.6)
  \entity{w}(0,-3.2){Wheel}
  \entity{b}(3.2,-3.2){Car body}
  \link{car}{p1}{1}\link{p1}{w}{n}
  \link{car}{p2}{1}\link{p2}{b}{1}
\end{erdiagram}
\end{example}

% =====================================================================
\section{Notes, subdomains and the grid}

\begin{center}
\verb|\note[|\meta{options}\verb|](|\meta{x,y}\verb|){|\meta{text}\verb|}|\qquad
\verb|\subdomain[|\meta{options}\verb|]{|\meta{label}\verb|}{(x,y) (x,y) ...}|\qquad
\verb|\grid(x1,y1)(x2,y2)|
\end{center}

Notes are free text (dark red by default; \texttt{color=} changes it; use
\verb|\\| for line breaks; TikZ options like \texttt{anchor=west} apply).
A subdomain is a frame drawn behind the diagram: two points give a
rectangle, more points a polygon; the label box hangs from the
\emph{first} point. Options: \texttt{color}, \texttt{dotted} (dotted
background), \texttt{fill=}\meta{colour} and TikZ options for the frame
(e.g.\ \texttt{rounded corners}, \texttt{dashed}). \cs{grid} draws a
labelled helper grid for a region (see also the environment option
\texttt{grid}).

For the package a subdomain is only a drawing. In the desktop editor it also groups
what is inside it: moving it moves its contents, and the relational schema is grouped under it.

\begin{example}
\begin{erdiagram}[max width=\linewidth]
  \subdomain[color=blue!60, dotted]{Sales}
    {(-1.6,0.9) (3,-1)}
  \subdomain[color=green!50!black, fill=green!5]
    {Logistics}{(3.4,0.9) (8,0.9) (8,-2.4)
    (5.5,-2.4) (5.5,-1) (3.4,-1)}
  \entity{c}(0,0){Customer}
  \entity{d}(6.5,0){Delivery}
  \note(1,-1.6){Customers must be\\at least 18}
  \note[color=blue, anchor=west](3.5,-2.1)
    {\textit{note in blue}}
\end{erdiagram}
\end{example}

% =====================================================================
\section{Relational model options}\label{sec:rm}

The desktop editor (\emph{ERDTeX}, menu \emph{File $\to$
Relational Schema / SQL}) transforms an ER diagram into relations and writes
them as plain text, as LaTeX, as Markdown, or as SQL for PostgreSQL, Oracle and SQLite. Where
the rules leave a choice, the choice is stored in the diagram as options. The
package accepts these options and ignores them, so such diagrams compile
unchanged.

Key attributes of a relationship type
(\verb|\attributes{r}{*date}|) join the primary key of its relation, and a role
(\verb|role=|) names the foreign key column. The two cardinality notations are described in section~\ref{sec:links}.

\begin{longtable}{@{}>{\ttfamily}p{4.4cm}p{11.4cm}@{}}
\toprule
\normalfont\textbf{Option} & \textbf{Meaning}\\\midrule
rm rules=db1 & in \cs{setup}: written by older versions of the editor (up to 3.5) for its
  second set of defaults (generalisations without categorisation, recursive
  relationships as foreign keys); still read by it\\
rm=1|2|3|4 & on \cs{isa}/\cs{canbe}: how the generalisation is represented: (1) one relation per entity type, (2) subtype
  relations only, (3) one relation with a categorisation attribute,
  (4) one relation with Boolean flags\\
rm category=\meta{kind} & with \texttt{rm=1}: \texttt{attribute}
  (categorisation attribute in the supertype, the default for disjoint
  generalisations), \texttt{relation} (own categorisation relation, the
  default for overlapping ones) or \texttt{none}\\
rm=\meta{entity id} & on a 1:1 or c:c relationship: the relation that
  receives the foreign key\\
rm=table & on any relationship that would be stored as foreign keys: an own
  relation instead\\
rm=fk & on a recursive relationship: a foreign key instead of an own
  relation\\
rm=fold & on an n-ary relationship: foreign keys in the relation of one side
  instead of an own relation\\
rm key=\meta{sides} & on n-ary relationships: the single sides that belong to
  the primary key, in addition to the complex ones (entity ids;
  \meta{id}\texttt{\#2} for the second link to the same entity). With (min,max):
  the side whose key identifies the relationship\\
rm name=\meta{name} & on entities and relationships: the name of the
  relation; on \cs{isa}: the name of the categorisation attribute\\
\bottomrule
\end{longtable}

The kinds of attributes have a fixed mapping. The parts of a composite attribute become
columns, with its name as a prefix only on a clash. A multi-valued attribute gets a relation of
its own, with the owner's key and the value. A derived attribute isn't stored; with a
\texttt{formula}, a view in the SQL script computes it. A further candidate key becomes
\texttt{UNIQUE}, where NULL stays allowed.

A subtype with several supertypes refers to all of them, and takes its key from the first.
A second supertype above the same root has the same key, so that column gets a second foreign
key; with different roots, the other key becomes a further column that is \texttt{UNIQUE} and
\texttt{NOT NULL}. The options \texttt{rm=2}, \texttt{3} and \texttt{4} aren't available for
such a generalisation, because the subtype couldn't be merged into two relations. A diagram
that asks for one gets a warning, and the default.

% =====================================================================
\subsection{Writing a relational schema}

The LaTeX output of the desktop editor is an environment of this package. It needs no other
package, and its commands only exist inside it, so a document can have its own
\cs{pkey} or \cs{set}. Primary keys are underlined, foreign keys in italics. Names stay
as drawn, only without spaces; SQL uses sanitised ones.

\begin{verbatim}
\begin{relationalschema}
  \relation{Raum}{\pkey{Raum-Id}{integer}, \col{Name}{string}}
  \relation{Sitz}{\pfkey{Raum-Id}{integer}, \pkey{Nr}{integer}, \fkey{Gast}{}}
\end{relationalschema}
\end{verbatim}
\begin{relationalschema}
  \relation{Raum}{\pkey{Raum-Id}{integer}, \col{Name}{string}}
  \relation{Sitz}{\pfkey{Raum-Id}{integer}, \pkey{Nr}{integer}, \fkey{Gast}{}}
\end{relationalschema}

\noindent Under \cs{relationgroup}\verb|{Title}| the desktop editor groups the relations of a
subdomain box.

\section{Settings for a whole diagram or document}\label{sec:settings}

Every style key can be set for one diagram (an option of the environment, or \cs{setup}
inside it) or for the rest of a document (\cs{ersetup}; in GSM diagrams \cs{gsmsetup}). A
diagram's own value wins. The desktop editor stores its settings per diagram with \cs{setup},
so they travel with the file.

\paragraph{Sans-serif.} \texttt{sans serif} sets all text of a diagram in the document's
sans-serif font, also in GSM diagrams. The text around the diagram stays as it is.

\paragraph{Scaling.} \texttt{scale=0.5} scales the finished diagram, with its nodes, text and
lines. It is an option of the environments and of \cs{inputerdiagram} and \cs{inputgsmdiagram},
just like \texttt{width}, \texttt{height} and \texttt{max width}. An exact \texttt{width} or
\texttt{height} wins over it; \texttt{max width} and \texttt{max height} still limit the result.

\paragraph{German labels.} The labels follow the language of the document (babel or polyglossia):
\emph{IS\_A}, \emph{CAN\_BE}, \emph{PART\_OF} become \emph{IST\_EIN*E}, \emph{KANN\_SEIN},
\emph{TEIL\_VON}. \texttt{label language=german} or \texttt{english} fixes the language; a label
of your own (\texttt{isa label=...}) always wins. The words of a generalisation's type
are not translated.

\begin{example}
\begin{erdiagram}[sans serif, label language=german, scale=0.8]
  \entity{p}(0,0){Person}
  \isa[constraint pos=east]{i}(0,-1.8){p}{cd}
  \entity{s}(0,-3.6){Student}
  \link{s}{i}
\end{erdiagram}
\end{example}

\section{Rule checks}\label{sec:rules}

While drawing, the package checks the most important ER rules and
reports violations as \LaTeX\ warnings (\texttt{Package erdiagram
Warning: ER rule: \dots}) without changing the output. They can be
switched off with \verb|\ersetup{rule warnings=false}| or, for the
end-of-diagram checks, with the environment option
\texttt{checks=false}.

\begin{itemize}
\item two entities linked directly, without a relationship in between;
\item two relationships linked directly;
\item a non-entity linked to an IS\_A, or an IS\_A whose supertype
  is not an entity;
\item a cardinality on an edge that touches no entity;
\item attributes attached to an IS\_A;
\item a relationship connected to fewer than two nodes;
\item an IS\_A without subtypes;
\item a weak entity that is not linked to a weak (E/Id) relationship;
\item a double line anywhere but between a weak relationship and a weak entity.
\end{itemize}

The desktop editor checks all of this and more, such as missing roles on recursive relationships
and invalid cardinalities (see its README).

% =====================================================================
\section{Style parameters}\label{sec:style}

All parameters can be set globally with \cs{ersetup}\verb|{key=value,...}|,
for one diagram as environment options, or inside a diagram with
\cs{setup}. Colours are xcolor expressions, fonts are font commands.

\begin{longtable}{@{}>{\ttfamily}p{5.1cm}>{\ttfamily}p{3.5cm}p{7cm}@{}}
\toprule
\normalfont\textbf{Key} & \normalfont\textbf{Default} & \textbf{Meaning}\\\midrule
\endhead
entity fill / draw / text & gray!15 / black / black & entity colours\\
entity line width & 0.6pt & \\
entity font & \textbackslash bfseries & \\
entity minimum width / height & 2.8cm / 1cm & \\
entity inner sep & 5pt & \\
weak double distance & 2pt & gap of double borders\\
rel fill / draw / text & white / black / black & relationship colours\\
rel line width & 0.6pt & \\
rel font & \textbackslash itshape & \\
rel minimum width / height & 2cm / 1.4cm & \\
rel aspect & 1.6 & width/height of diamonds\\
dependency color / font & red!70!black / \textbackslash normalfont\textbackslash bfseries & the E / Id mark\\
isa label / canbe label & IS\textbackslash\_A / CAN\textbackslash\_BE & \\
partof label & PART\textbackslash\_OF & \\
constraint color / font & red!70!black / \textbackslash normalfont\textbackslash small & cd, co, id, io\\
constraint pos & inside & see \cs{isa}\\
constraint style & letters & or \texttt{words}\\
isa shape & hexagon & or \texttt{diamond}\\
isa minimum width / height & 2cm / 1cm & of the hexagon\\
isa slant & 0.5 & how pointed the hexagon's sides are\\
attr font / text & \textbackslash small / black & attribute trees\\
attr spread & 0.35cm & distance between attributes\\
attr indent & 3pt & the parts of a composite attribute\\
attribute style & compact & or \texttt{chen}: lists as ovals\\
sans serif & false & section~\ref{sec:settings}\\
label language & auto & or \texttt{german}, \texttt{english}\\
chen dist & 0.9cm & from a node to its ovals\\
show types & false & print data types (section~\ref{sec:types})\\
type sep / type font & \verb*|:\ | / (empty) & between name and type / font of the type\\
attr gap & (= spread) & node to first attribute (trunk)\\
attr tick / attr stem & 0.3cm / 0.3cm & \\
attribute fill / draw & white / black & Chen-style ellipses\\
attribute minimum width / height & 1.2cm / 0.7cm & \\
card font / text & \textbackslash footnotesize / black & cardinalities\\
card placement & inline & or \texttt{beside}\\
card symbols & labels & or \texttt{crowsfoot}\\
crow length / width & 6pt / 3.5pt & the crow's foot symbols\\
card notation & n & or \texttt{minmax}; what the desktop editor offers\\
card dist & 0.3cm & \\
card side & auto & \\
role font / text & \textbackslash footnotesize \textbackslash itshape / black & role names\\
role pos & auto & 0.5, or near the entity if the cardinality is mid-edge\\
edge color / line width & black / 0.5pt & \\
link double distance & 1.5pt & gap of double-lined links\\
note color / font & red!70!black / \textbackslash small & \\
subdomain color & blue!60 & \\
subdomain label font & \textbackslash small\textbackslash bfseries & \\
subdomain line width & 1.2pt & \\
subdomain tuck & 4pt & chamfer of the label box\\
subdomain label vpad / hpad & 4pt / 5pt & \\
rule warnings & true & \\
grid color / grid label color & gray!25 / gray!55 & \\
\bottomrule
\end{longtable}

The TikZ styles \texttt{er entity}, \texttt{er relationship},
\texttt{er weak}, \texttt{er attribute}, \texttt{er attr label},
\texttt{er card}, \texttt{er role}, \texttt{er constraint},
\texttt{er edge}, \texttt{er arrow} and \texttt{er note} can be
redefined or extended with \cs{tikzset} for deeper changes.

% =====================================================================
\section{A complete example: a music festival}

The sample diagram of ERDTeX (\emph{File $\to$ Open Example}),
drawn with the package. It uses most features at once: a weak entity type
with its identifying relationship, generalisations of both kinds, PART\_OF,
attributes of every kind with data types, a recursive and a three-way
relationship (through a junction point), subdomain boxes and a note.

The diagram is stored in the file \texttt{festival.erd} (listed below) and
included like this:

\begin{center}
\verb|\inputerdiagram[width=\linewidth]{festival.erd}|
\end{center}

\medskip
\noindent\inputerdiagram[width=\linewidth]{festival.erd}

\lstinputlisting[basicstyle=\ttfamily\scriptsize, frame=single,
  rulecolor=\color{black!25}]{festival.erd}

% =====================================================================
\section{Known limitations}\label{sec:limits}

The following cases are \emph{not} handled automatically (yet). They are
shown here on purpose.

\begin{enumerate}
\item \textbf{Labels only look at directions.} A cardinality chooses the
  side of its edge by comparing the \emph{directions} of the other edges
  and attribute trees at its entity, and it keeps clear of the entity
  itself, but it does not test for collisions with the actual drawings
  of other elements. Below, the edge and the attribute tree leave A in
  the same direction, so the label collides with the tree. Nothing
  prevents nodes or attribute trees of different elements from
  overlapping either: positions are up to the author (or the editor).
\begin{example}
\begin{erdiagram}
  \entity{a}(0,0){A}
  \relationship{r}(3.5,0){R}
  \entity{b}(3.5,-2){B}
  \link{a}{r}{(0,n)}
  \link{r}{b}{(1,1)}
  \attributes[anchor=east]{a}{x, y}
\end{erdiagram}
\end{example}
\item \textbf{Very crowded entities.} When three or more edges leave an
  entity within a narrow angle, some cardinality ends up between two
  edges and may touch one of them; use \texttt{card side} or
  \texttt{card dist} for such edges.
\item \textbf{Way points} in \texttt{via} must be plain coordinates
  \verb|(x,y)| or node anchors \verb|(a.north)|; calculations like
  \verb|($(a)+(1,0)$)| are not supported there (use a \cs{point}).
\item \textbf{Polygon subdomains} hang their label box from the first
  point, so list the top-left corner first.
\item \textbf{Rule checks} are local: e.g.\ cycles in generalisation
  hierarchies are not detected.
\end{enumerate}

\end{document}
